\begin{center}
{\large\bfseries Abstract}
\end{center}

\todo{write after results; this is a placeholder of the right length and shape}

The healthy heartbeat is neither regular nor random. Its beat-to-beat fluctuations are correlated across time scales, and the short-term scaling exponent of detrended fluctuation analysis (\aone) summarises this organisation in a single number that remains informative during physical activity. Movement practices built around handling an object are often described by practitioners as absorbing and restorative, but whether they leave a measurable acute trace in autonomic regulation has not been tested.

The present study recorded RR intervals with a Polar H10 chest strap while participants rested and while they performed two rhythmic stick-manipulation movements, isolations and spins, each for two minutes. \al{} was estimated per trial with an open analysis pipeline and compared between rest and movement, between the two movements, and between participants of differing expertise, with mean heart rate as a covariate to separate effects of skill from effects of physical load.

\todo{results in two to three sentences: direction of the shift, whether it survives the heart-rate adjustment, movement and expertise effects}

\todo{one to two sentences of interpretation and what follows for the stick as a research tool}

\textbf{Keywords:} heart rate variability, detrended fluctuation analysis, DFA alpha1, fractal physiology, object manipulation, motor skill, expertise, flow, Polar H10
